\section{Fisher}

\begin{definition}[Fisher 精准性检验]
    Fisher 精准性检验是一种用于检验两个变量之间的相关关系的方法，与 $K^2$ 检验有很多相似之处，但是它仅适用于小样本的校验，一般在 $K^2$ 检验失效时（样本过小）使用。

    Fisher 精准性检验通过计算精确的 $p$ 值来判断两个分类变量之间的关联性，适用于 $2 \times 2$ 列联表。其公式如下所示：

    \begin{center}
        \begin{tabular}{|c|c|c|c|}
            \hline                  & Men              & Women            & 女性            \\
            \hline Studying 学习      & $\boldsymbol{a}$ & $\boldsymbol{b}$ & $a+b$         \\
            \hline Non－studying 非学习 & $\boldsymbol{c}$ & $\boldsymbol{d}$ & $c+d$         \\
            \hline 列总计              & $a+c$            & $b+d$            & $a+b+c+d(=n)$ \\
            \hline
        \end{tabular}
    \end{center}

    Fisher showed that conditional on the margins of the table，$a$ is distributed as a hypergeometric distribution with $a+c$ draws from a population with $a+b$ successes and $c+d$ failures．The probability of obtaining such set of values is given by：

    Fisher 证明，在表格边界条件下， $\boldsymbol{a}$ 服从超几何分布，其中 $\boldsymbol{a}+\boldsymbol{c}$ 取自于 $\boldsymbol{a}+\boldsymbol{b}$ 个成功案例和 $\boldsymbol{c}+\boldsymbol{d}$ 个失败案例的总体。获得此类值的概率如下：
    $$
        p=\frac{\binom{a+b}{a}\binom{c+d}{c}}{\binom{n}{a+c}}=\frac{\binom{a+b}{b}\binom{c+d}{d}}{\binom{n}{b+d}}=\frac{(a+b)!(c+d)!(a+c)!(b+d)!}{a!b!c!d!n!}
    $$
\end{definition}

\begin{table}[H]
    \centering
    \caption{卡方检验与 Fisher 精准检验的对比}
    \label{tab:chi2_vs_fisher}
    % l: 第一列左对齐
    % X: 后两列为自适应宽度的段落列，平分剩余空间
    \begin{tabularx}{\textwidth}{l c c}
        \toprule
        \textbf{特性}                          & \textbf{卡方检验 ($\chi^2$)} & \textbf{Fisher 精准检验} \\
        \midrule

        \textbf{检验类型}                        &
        近似检验                                 &
        精准检验                                                                                   \\

        \addlinespace % 在行之间增加一点额外空间，让表格更疏朗

        \textbf{适用条件}                        &
        \textbf{大样本数据}。至少 80\% 的单元格的期望频数大于 5 &
        \textbf{小样本数据}。当卡方检验的条件不满足时使用                                                          \\
        \bottomrule
    \end{tabularx}
\end{table}

\begin{table}[H]
    \centering
    \caption{新药疗效检验的列联表}
    \label{tab:drug_trial}
    \begin{tabular}{l c c c}
        \toprule
        \textbf{组别} & \textbf{康复} & \textbf{未康复} & \textbf{总计} \\
        \midrule
        新药组         & 8           & 2            & \textbf{10} \\
        安慰剂组        & 3           & 7            & \textbf{10} \\
        \midrule
        \textbf{总计} & \textbf{11} & \textbf{9}   & \textbf{20} \\
        \bottomrule
    \end{tabular}
\end{table}

在表 \ref{tab:drug_trial} 中，我们展示了新药组与安慰剂组的康复情况对比。

\begin{lstlisting}[style=python]
import numpy as np
from scipy.stats import fisher_exact

# 2 x 2 列联表
# [8 , 2, 10]
# [3 , 7, 10]
# [11, 9, 20]
table = np.array([[8, 2], [3, 7]])
odds_ratio, p_fisher = fisher_exact(table) # fisher_exact 会返回: 优势比 (Odds Ratio), p值
print(f"优势比: {odds_ratio:.4f} P 值: {p_fisher:.4f}")

alpha = 0.05
if p_fisher < alpha:
    print(f"在显著性水平 {alpha} 下，P-value ({p_fisher:.4f}) < {alpha}，我们拒绝原假设。")
    print("能够认为两个变量之间存在统计学关联性。")
else:
    print(f"在显著性水平 {alpha} 下，P-value ({p_fisher:.4f}) >= {alpha}，我们不能拒绝原假设。")
    print("无法证明两个变量之间有关联。")
\end{lstlisting}